% Book 3, Year 1, chapter 27 -- Boolean algebra and combinational circuits.
% Provenance (never printed): sources/book3/BRIEFS.md, chapter 27.
% ledger: b1:boole-1854, b1:shannon-1938, b1:karnaugh-1953, b1:sheffer-1913
\chapter{Boolean Algebra and Combinational Circuits}\label{ch:b1:combinational-circuits}

The adder inside a processor adds two 64-bit numbers in a fraction of a nanosecond. It is built from gates, each of
which knows only one or two bits and computes a function as simple as ``both'' or ``either''. This chapter is about how
such gates are combined: how any function of bits can be written with them, how to write it with few, why one kind of
gate is enough, and how the arrangement of the gates, not their number, decides the speed of an adder. The circuits
are built and run in a small simulator, in which a circuit is what it is in the hardware: a \omterm{def:b1:graph-traversals:dag}{directed acyclic graph} of
gates.

\begin{recall}
Gates and the half adder were built in the first book; the second built an adder and a \omterm{def:b1:combinational-circuits:mux}{multiplexer}. Formulas, their
truth tables and the normal forms DNF and CNF are in \cref{def:b1:propositional-logic:normal-forms}; DAGs and their
longest paths in \cref{ch:b1:graph-traversals,ch:b1:shortest-paths}.
\end{recall}

\section{Boolean functions and canonical forms}\label{sec:b1:combinational-circuits:functions}

\begin{definition}[Boolean function]\label{def:b1:combinational-circuits:boolean-function}
A \emph{Boolean function}\index{Boolean function} of $n$ variables is a function $f\colon \{0, 1\}^n \to \{0, 1\}$. It is
given by its truth table, $2^n$ rows; so there are $2^{2^n}$ Boolean functions of $n$ variables.
\end{definition}

We write $x'$ for the negation of $x$, $xy$ for the conjunction and $x + y$ for the disjunction. These operations obey
the laws of \emph{Boolean algebra}\index{Boolean algebra}: commutativity and associativity of $+$ and of the product;
each distributes over the other, $x(y + z) = xy + xz$ and $x + yz = (x + y)(x + z)$; $x + 0 = x$, $x \cdot 1 = x$, $x + x'
= 1$, $x x' = 0$; and De Morgan's laws $(x + y)' = x' y'$, $(xy)' = x' + y'$. Each is checked on the four (or eight)
rows of a truth table. Boole set them out in 1854 as laws of thought; in 1938 Shannon showed that they describe
circuits of relays.

\omimg{book3/george-boole.jpg}{0.22\linewidth}{George Boole, whose book of 1854 set out these laws as the laws of
thought. Portrait by an unknown artist, public domain.}

\begin{definition}[Minterm, maxterm, canonical forms]\label{def:b1:combinational-circuits:canonical}
A \emph{minterm}\index{minterm} of $n$ variables is a product in which every variable appears once, plain or negated: it
is 1 on exactly one row. A \emph{maxterm}\index{maxterm} is a sum in which every variable appears once: it is 0 on exactly
one row. The \emph{canonical sum of products}\index{canonical sum of products} of $f$ is the sum of the minterms of the
rows where $f$ is 1; its \emph{canonical product of sums}\index{canonical product of sums} is the product of the maxterms
of the rows where $f$ is 0.
\end{definition}

Both equal $f$: the sum is 1 exactly on the rows of its \omterm{def:b1:combinational-circuits:canonical}{minterms}, and the product 0 exactly on the rows of its
\omterm{def:b1:combinational-circuits:canonical}{maxterms}. They are the DNF and CNF of \cref{ch:b1:propositional-logic} read off a truth table, and they prove that
negation, conjunction and disjunction can express every \omterm{def:b1:combinational-circuits:boolean-function}{Boolean function}.

\omcode{code/bachelor-1/27-combinational-circuits/listings/canonical.py}{15}{29}{The canonical forms from the truth
table}\label{lst:b1:combinational-circuits:canonical}

\omoutput{code/bachelor-1/27-combinational-circuits/outputs/canonical.txt}{The majority of three bits, its two
canonical forms, and the number of functions of 1 to 5 variables.}\label{out:b1:combinational-circuits:canonical}

\begin{definition}[Combinational circuit]\label{def:b1:combinational-circuits:combinational}
A \emph{combinational circuit}\index{combinational circuit} is a \omterm{def:b1:graph-traversals:dag}{directed acyclic graph} whose sources are the inputs,
whose other \omterm{def:b1:graph-traversals:graph}{vertices} are logic gates, each fed by the \omterm{def:b1:graph-traversals:graph}{vertices} it has \omterm{def:b1:graph-traversals:graph}{arcs} from, and some of whose \omterm{def:b1:graph-traversals:graph}{vertices} are marked as
outputs. Its outputs depend only on its current inputs.
\end{definition}

The simulator stores the gates in the order they were added; a gate may only use wires that exist already, so that
order is topological, and evaluation is one pass, as for the DAG paths of \cref{ch:b1:shortest-paths}.

\omcode{code/bachelor-1/27-combinational-circuits/listings/circuit.py}{20}{40}{The simulator: adding a gate, evaluating,
and the delay}\label{lst:b1:combinational-circuits:circuit}

\section{Simplification}\label{sec:b1:combinational-circuits:simplification}

The majority function needs four products of three literals in its canonical form, but $ab + ac + bc$ will do. Algebra
finds it: $abc$ can be repeated, $a'bc + abc = bc$, and so on. A picture finds it faster.

\begin{definition}[Karnaugh map]\label{def:b1:combinational-circuits:karnaugh}
A \emph{Karnaugh map}\index{Karnaugh map} of a function of four variables is its truth table arranged as a $4 \times 4$
grid, rows labelled by $ab$ and columns by $cd$ in the order $00, 01, 11, 10$, so that two neighbouring cells, the
first and last of a row or column included, differ in exactly one variable.
\end{definition}

Two neighbouring 1 cells merge into a product without the variable in which they differ; four cells in a $2 \times 2$
square or a row of four lose two variables; eight cells lose three. Groups must be rectangles of $1$, $2$, $4$ or $8$
cells, possibly wrapping around the edges.

\begin{definition}[Implicant, prime implicant, essential prime implicant, don't-care]\label{def:b1:combinational-circuits:implicant}
An \emph{implicant}\index{implicant} of $f$ is a product of literals that is 1 only where $f$ is 1. It is
\emph{prime}\index{prime implicant} if removing any literal makes it cover a 0 of $f$, and \emph{essential}\index{essential
prime implicant} if it is the only prime implicant covering some 1 of $f$. A \emph{don't-care
condition}\index{don't-care condition} is a row whose output does not matter, because it never occurs; it may be
covered or not, as suits the simplification.
\end{definition}

A minimal sum of products is made of \omterm{def:b1:combinational-circuits:implicant}{prime implicants} (a \omterm{def:b1:combinational-circuits:implicant}{non-prime} one can lose a literal) and contains every \omterm{def:b1:combinational-circuits:implicant}{essential}
one. The program finds the \omterm{def:b1:combinational-circuits:implicant}{prime implicants} by merging terms that differ in one variable, round after round, the method
of Quine and McCluskey; with four variables, a \omterm{def:b1:combinational-circuits:karnaugh}{Karnaugh map} shows the same groups at a glance.

\omcode{code/bachelor-1/27-combinational-circuits/listings/karnaugh.py}{19}{40}{\omterm{def:b1:combinational-circuits:implicant}{Prime implicants} by repeated merging;
\code{-} marks a variable that has gone}\label{lst:b1:combinational-circuits:qm}

\omoutput{code/bachelor-1/27-combinational-circuits/outputs/karnaugh.txt}{The map of a function of four variables, its six
\omterm{def:b1:combinational-circuits:implicant}{prime implicants} and the two \omterm{def:b1:combinational-circuits:implicant}{essential} ones.}\label{out:b1:combinational-circuits:qm}

\begin{omfigure}
\begin{tikzpicture}[x=9mm, y=-7mm, font=\small]
\foreach \lab [count=\j from 0] in {00,01,11,10} {\node at (\j,-1) {\lab}; \node at (-1,\j) {\lab};}
\node at (-1,-1) {\scriptsize $ab \backslash cd$};
\foreach \i/\row in {0/{1,1,0,1}, 1/{0,1,1,1}, 2/{0,0,0,1}, 3/{1,1,0,1}}
  \foreach \v [count=\j from 0] in \row {\draw[gray] (\j-0.5,\i-0.5) rectangle (\j+0.5,\i+0.5);
    \ifnum\v=1 \node at (\j,\i) {1};\fi}
\draw[omDef, very thick, rounded corners=2mm] (2.65,-0.4) rectangle (3.35,3.4);
\draw[omThm, very thick, rounded corners=2mm] (-0.4,-0.65) -- (-0.4,0.35) -- (1.4,0.35) -- (1.4,-0.65);
\draw[omThm, very thick, rounded corners=2mm] (-0.4,3.65) -- (-0.4,2.65) -- (1.4,2.65) -- (1.4,3.65);
\draw[black!70, very thick, dashed, rounded corners=2mm] (0.62,0.62) rectangle (2.38,1.38);
\node[anchor=west, align=left, font=\scriptsize] at (4,1.5) {\textcolor{omDef}{$cd'$: the column $10$}\\
\textcolor{omThm}{$b'c'$: rows $00$ and $10$ wrap round}\\ $a'bd$ (dashed): covers 5 and 7};
\end{tikzpicture}
\par\smallskip{\small The map of \cref{out:b1:combinational-circuits:qm} covered by three \omterm{def:b1:combinational-circuits:implicant}{prime implicants}: $f = cd' +
b'c' + a'bd$. The two \omterm{def:b1:combinational-circuits:implicant}{essential} ones leave the cells 5 and 7 bare; $a'bd$ covers both, where $a'c'd$ or $a'bc$ would
cover only one.}
\end{omfigure}

A seven-segment display shows a digit with seven lit bars, named $a$ to $g$. Its \omterm{def:b1:combinational-circuits:mux}{decoder} takes the digit in four bits;
the codes 10 to 15 never occur, and the don't-cares let the groups grow.

\omoutput{code/bachelor-1/27-combinational-circuits/outputs/seven_segment.txt}{The seven segments, each simplified with
the codes 10 to 15 as don't-cares, and checked on the ten digits. The variables are $abcd$, $a$ the high
bit.}\label{out:b1:combinational-circuits:seven}

\section{NAND alone is enough}\label{sec:b1:combinational-circuits:nand}

\begin{definition}[Functionally complete set]\label{def:b1:combinational-circuits:complete}
A set of Boolean operations is \emph{functionally complete}\index{functionally complete set} if every \omterm{def:b1:combinational-circuits:boolean-function}{Boolean function} of
any number of variables can be computed by a circuit of gates from the set.
\end{definition}

The canonical forms show that $\{\text{NOT}, \text{AND}, \text{OR}\}$ is complete; by De Morgan, so are $\{\text{NOT},
\text{AND}\}$ and $\{\text{NOT}, \text{OR}\}$.

\begin{proposition}[NAND is complete]\label{prop:b1:combinational-circuits:nand}
$\{\text{NAND}\}$ is \omterm{def:b1:combinational-circuits:complete}{functionally complete}.
\end{proposition}

\begin{proof}
$x' = \text{NAND}(x, x)$ and $xy = \text{NAND}(x, y)'$, so NOT and AND are expressed with NAND, and $\{\text{NOT},
\text{AND}\}$ is complete.
\end{proof}

\omcode{code/bachelor-1/27-combinational-circuits/listings/nand_only.py}{9}{23}{NOT, AND, OR and XOR from NAND
gates}\label{lst:b1:combinational-circuits:nand}

\omoutput{code/bachelor-1/27-combinational-circuits/outputs/nand_only.txt}{Each construction checked on its four
inputs.}\label{out:b1:combinational-circuits:nand}

\begin{omfigure}
\begin{circuitikz}[font=\small]
\node[nand port] (n) at (1.6,0) {};
\coordinate (j) at ($(n.in 1)!0.5!(n.in 2)+(-0.35,0)$);
\draw (j) -- ++(-0.5,0) node[left] {$x$}; \draw (j) |- (n.in 1); \draw (j) |- (n.in 2); \fill (j) circle (1.3pt);
\draw (n.out) -- ++(0.3,0) node[right] {$x'$};
\node at (1.4,-1.1) {NOT};
\begin{scope}[xshift=43mm]
\node[nand port] (a1) at (1.0,0) {}; \node[nand port] (a2) at (3.2,0) {};
\draw (a1.in 1) -- ++(-0.4,0) node[left] {$x$}; \draw (a1.in 2) -- ++(-0.4,0) node[left] {$y$};
\coordinate (k) at ($(a2.in 1)!0.5!(a2.in 2)+(-0.35,0)$);
\draw (a1.out) -- (k); \draw (k) |- (a2.in 1); \draw (k) |- (a2.in 2); \fill (k) circle (1.3pt);
\draw (a2.out) -- ++(0.3,0) node[right] {$xy$}; \node at (2.1,-1.1) {AND};
\end{scope}
\begin{scope}[xshift=96mm]
\node[nand port] (o1) at (1.2,0.75) {}; \node[nand port] (o2) at (1.2,-0.75) {}; \node[nand port] (o3) at (3.2,0) {};
\coordinate (p) at ($(o1.in 1)!0.5!(o1.in 2)+(-0.35,0)$); \coordinate (q) at ($(o2.in 1)!0.5!(o2.in 2)+(-0.35,0)$);
\draw (p) -- ++(-0.4,0) node[left] {$x$}; \draw (p) |- (o1.in 1); \draw (p) |- (o1.in 2);
\draw (q) -- ++(-0.4,0) node[left] {$y$}; \draw (q) |- (o2.in 1); \draw (q) |- (o2.in 2);
\fill (p) circle (1.3pt) (q) circle (1.3pt);
\draw (o1.out) -- ++(0.25,0) |- (o3.in 1); \draw (o2.out) -- ++(0.25,0) |- (o3.in 2);
\draw (o3.out) -- ++(0.3,0) node[right] {$x + y$};
\node at (2.0,-1.75) {OR, by De Morgan: $x + y = (x' y')'$};
\end{scope}
\end{circuitikz}
\par\smallskip{\small NOT, AND and OR from NAND gates, as in \cref{lst:b1:combinational-circuits:nand}: one, two and
three gates.}
\end{omfigure}

Not every set is complete. A function is \emph{monotone} if raising an input from 0 to 1 never lowers the output; AND
and OR are monotone, and so is any circuit of AND and OR gates, by induction along the DAG. NOT is not monotone, so
$\{\text{AND}, \text{OR}\}$ is not complete (\cref{exo:b1:combinational-circuits:10}). Sheffer showed in 1913 that one
operation, the ``stroke'' now called NAND, suffices for all of Boolean algebra. In hardware a NAND gate is the
cheapest: a few transistors used as switches, the physics of which belongs to the physics books.

\begin{history}[Gates on a chip]
% ledger: b1:7400, b1:7400-gates
In the 1960s, gates stopped being wired one by one from separate parts and began to be made together on one chip of
silicon, an integrated circuit. In 1966 Texas Instruments announced its 7400 series of such chips in cheap plastic
packages. The 7400 itself holds four independent two-input NAND gates in a package of 14 pins: three for each gate,
and two for the power supply. Circuits like those of this chapter were then built by wiring such chips together on a
board.
\end{history}

\omimg{book3/sn7400n-nand-chip.jpg}{0.4\linewidth}{A 7400 chip in its plastic package with two rows of seven pins.
Inside are the four NAND gates; the notch marks the end where pin 1 is. Photograph: Mister rf, CC BY-SA 4.0.}

\section{Arithmetic and routing}\label{sec:b1:combinational-circuits:arithmetic}

\begin{definition}[Full adder, ripple-carry adder]\label{def:b1:combinational-circuits:adder}
A \emph{full adder}\index{full adder} adds three bits $a$, $b$ and a carry in $c$: its sum is $a \oplus b \oplus c$ and its
carry out is $ab + c(a \oplus b)$, where $\oplus$ is XOR. A \emph{ripple-carry adder}\index{ripple-carry adder} of $n$
bits chains $n$ full adders, the carry out of each feeding the carry in of the next.
\end{definition}

Subtraction comes free: in two's complement, $a - b = a + b' + 1$, where $b'$ inverts every bit; XOR gates with a control
bit invert $b$ on demand, and the same bit is the first carry in.

The carry is the trouble: bit $i$ cannot finish before bit $i - 1$, and the last carry waits for the whole chain. The way
out is to compute every carry directly from the inputs. Write $g_i = a_i b_i$ (bit $i$ \emph{generates} a carry) and
$p_i = a_i \oplus b_i$ (it \emph{propagates} the carry it receives). Then $c_{i+1} = g_i + p_i c_i$, and unrolling,
\[ c_{i+1} = g_i + p_i g_{i-1} + p_i p_{i-1} g_{i-2} + \dots + p_i p_{i-1} \cdots p_0 c_0. \]

\begin{definition}[Generate and propagate, carry-lookahead adder]\label{def:b1:combinational-circuits:lookahead}
The signals $g_i$ and $p_i$ are the \emph{generate and propagate signals}\index{generate and propagate signals}. A
\emph{carry-lookahead adder}\index{carry-lookahead adder} computes the carries of a block of bits by the unrolled formula,
with one level of AND gates and one OR gate per carry.
\end{definition}

\omcode{code/bachelor-1/27-combinational-circuits/listings/adders.py}{23}{35}{A 4-bit lookahead
block}\label{lst:b1:combinational-circuits:lookahead}

\begin{omfigure}
\begin{tikzpicture}[b/.style={draw, minimum width=11mm, minimum height=9mm, font=\small}, >=stealth, font=\scriptsize]
\foreach \i in {0,...,3} {\node[b] (fa\i) at ({-1.5*\i},0) {FA}; \node at ({-1.5*\i},0.85) {$a_\i\ b_\i$};
  \draw[->] ({-1.5*\i},0.65) -- (fa\i.north); \draw[->] (fa\i.south) -- ++(0,-0.4) node[below] {$s_\i$};}
\foreach \i/\j in {0/1, 1/2, 2/3} \draw[->, omThm, thick] (fa\i.west) -- node[above] {$c_\j$} (fa\j.east);
\draw[<-, omThm, thick] (fa0.east) -- ++(0.6,0) node[right] {$c_0$};
\draw[->, omThm, thick] (fa3.west) -- ++(-0.6,0) node[left] {$c_4$};
\node at (-2.25,-1.6) {ripple: the carry crosses every stage};
\begin{scope}[xshift=38mm]
\node[b, minimum width=34mm] (gp) at (1.6,0.9) {$g_i = a_i b_i$, $p_i = a_i \oplus b_i$};
\node[b, minimum width=34mm, fill=omDef!12] (cl) at (1.6,-0.5) {carries $c_1 \dots c_4$ at once};
\node[b, minimum width=34mm] (s) at (1.6,-1.9) {$s_i = p_i \oplus c_i$};
\draw[->] (gp) -- (cl); \draw[->] (cl) -- (s); \draw[<-] (cl.east) -- ++(0.5,0) node[right] {$c_0$};
\draw[<-] (gp.north) -- ++(0,0.35) node[above] {$a_3 \dots a_0$, $b_3 \dots b_0$};
\node[below=1mm] at (s.south) {4 gate delays};
\end{scope}
\end{tikzpicture}
\par\smallskip{\small Left, a 4-bit \omterm{def:b1:combinational-circuits:adder}{ripple-carry adder}: each \omterm{def:b1:combinational-circuits:adder}{full adder} waits for the carry of the one before. Right, a
lookahead block: generate and propagate in one gate delay, all four carries in two more, the sums in one more.}
\end{omfigure}

\begin{definition}[Multiplexer, decoder]\label{def:b1:combinational-circuits:mux}
A $2^k$-to-1 \emph{multiplexer}\index{multiplexer} has $2^k$ data inputs and $k$ select inputs, and outputs the data input
whose number the select inputs spell. A $k$-to-$2^k$ \emph{decoder}\index{decoder} has $k$ inputs and $2^k$ outputs, of
which exactly the one numbered by the inputs is 1.
\end{definition}

A 2-to-1 \omterm{def:b1:combinational-circuits:mux}{multiplexer} is $s' x_0 + s x_1$; a \omterm{def:b1:combinational-circuits:mux}{decoder} is one AND gate per \omterm{def:b1:combinational-circuits:canonical}{minterm}. \omterm{def:b1:combinational-circuits:mux}{Multiplexers} \emph{route} information:
they choose which result goes on.

\begin{definition}[Arithmetic and logic unit]\label{def:b1:combinational-circuits:alu}
An \emph{arithmetic and logic unit}\index{arithmetic and logic unit} (ALU) computes, from two operands and an operation
code, one of several results, chosen by \omterm{def:b1:combinational-circuits:mux}{multiplexers}, and \emph{flags} that describe it: $Z$ (zero), $N$ (negative, the
top bit), $C$ (carry out) and $V$ (overflow in two's complement).
\end{definition}

\omcode{code/bachelor-1/27-combinational-circuits/listings/alu.py}{15}{33}{A 4-bit ALU: AND, OR, addition and
subtraction}\label{lst:b1:combinational-circuits:alu}

\omoutput{code/bachelor-1/27-combinational-circuits/outputs/alu.txt}{Checked against integer arithmetic on all 1024
cases; $7 + 1$ overflows in signed 4-bit arithmetic, and $9 + 9$ carries out.}\label{out:b1:combinational-circuits:alu}

Overflow happens when the two operands of the addition have the same sign and the result the other one; the circuit
tests exactly that, on the top bits.

\section{Propagation delay}\label{sec:b1:combinational-circuits:delay}

\begin{definition}[Propagation delay, longest delay path]\label{def:b1:combinational-circuits:delay}
Each gate takes a small time, its \emph{propagation delay}\index{propagation delay}, to update its output. The delay of a
circuit is the largest total delay along a path from an input to an output, a \emph{longest delay path}\index{longest
delay path}. With every gate counted as one delay, as in the simulator, it is the longest path of the DAG, counted in
gates.
\end{definition}

\omoutput{code/bachelor-1/27-combinational-circuits/outputs/adders.txt}{Sizes and delays of the two adders, both checked
on every input for 4 and 8 bits.}\label{out:b1:combinational-circuits:adders}

The \omterm{def:b1:combinational-circuits:adder}{ripple-carry adder} has delay $2n + 1$ (\cref{exo:b1:combinational-circuits:11}): linear in the number of bits. A
lookahead block has delay 4 whatever the carry in, at the price of more gates, and of gates with many inputs, which real
hardware limits. Rippling between blocks of four still costs two delays per block. Lookahead applied again to the blocks,
in a tree of lookahead units, brings the delay to $O(\log n)$; processors add 64 bits that way, within a clock cycle.

\section{Exercises}\label{sec:b1:combinational-circuits:exercises}

\begin{exercise}[$\star$]\label{exo:b1:combinational-circuits:1}
Give the truth table of the circuit $f = \text{OR}(\text{AND}(a, \text{NOT}(b)), \text{AND}(b, c))$, and its delay.
\end{exercise}

\begin{exercise}[$\star$]\label{exo:b1:combinational-circuits:2}
The function $f(a, b, c)$ is 1 exactly when an odd number of its inputs are 1. Write its \omterm{def:b1:combinational-circuits:canonical}{canonical sum of products}, and
explain why no Karnaugh grouping shortens it.
\end{exercise}

\begin{exercise}[$\star$]\label{exo:b1:combinational-circuits:3}
Using the map of \cref{out:b1:combinational-circuits:qm}, find another minimal sum of products, if there is one, and say
why every minimal sum has at least three products.
\end{exercise}

\begin{exercise}[$\star$]\label{exo:b1:combinational-circuits:4}
What does the ALU of \cref{lst:b1:combinational-circuits:alu} output for the operation code 11 with $a = 3$ and $b = 9$?
Give the four flags, and interpret them as unsigned and as signed numbers.
\end{exercise}

\begin{exercise}[$\star\star$]\label{exo:b1:combinational-circuits:5}
A student labels the columns of a \omterm{def:b1:combinational-circuits:karnaugh}{Karnaugh map} $00, 01, 10, 11$ and the rows likewise, then groups the two middle 1s of
the row $ab = 00$ of a function that is 1 on the \omterm{def:b1:combinational-circuits:canonical}{minterms} 1 and 2 only, writing $a'b'$ times ``whatever $cd$ is''. What
went wrong, and what is the correct sum?
\end{exercise}

\begin{exercise}[$\star\star$]\label{exo:b1:combinational-circuits:6}
This \omterm{def:b1:combinational-circuits:adder}{full adder} assembles its carry differently. On which row is it wrong, and why on that row only?
\omcode{code/bachelor-1/27-combinational-circuits/listings/exo_carry_xor.py}{7}{10}{Another carry}
\omoutput[2-2]{code/bachelor-1/27-combinational-circuits/outputs/exo_carry_xor.txt}{Every row checked.}
\end{exercise}

\begin{exercise}[$\star\star$]\label{exo:b1:combinational-circuits:7}
Build a 2-to-1 and a 4-to-1 \omterm{def:b1:combinational-circuits:mux}{multiplexer} in the simulator, and test them on every input.
\end{exercise}

\begin{exercise}[$\star\star$]\label{exo:b1:combinational-circuits:8}
Build a circuit that outputs $a < b$ for two 4-bit unsigned numbers, test it on every input, and give its delay.
\end{exercise}

\begin{exercise}[$\star\star$]\label{exo:b1:combinational-circuits:9}
Count the monotone \omterm{def:b1:combinational-circuits:boolean-function}{Boolean functions} of $n$ variables for $n \leq 3$, by program and by argument for $n \leq 2$.
\end{exercise}

\begin{exercise}[$\star\star\star$]\label{exo:b1:combinational-circuits:10}
Prove that $\{\text{NOR}\}$ is \omterm{def:b1:combinational-circuits:complete}{functionally complete}, and that $\{\text{AND}, \text{OR}\}$ is not.
\end{exercise}

\begin{exercise}[$\star\star\star$]\label{exo:b1:combinational-circuits:11}
Prove that in the simulator's model the $n$-bit \omterm{def:b1:combinational-circuits:adder}{ripple-carry adder} has delay $2n + 1$ for $n \geq 1$, and the lookahead
block delay 4.
\end{exercise}

\begin{exercise}[$\star\star\star$]\label{exo:b1:combinational-circuits:12}
Build a 4-bit \omterm{def:b1:pointers-and-arrays:array}{array} multiplier from AND gates and \omterm{def:b1:combinational-circuits:adder}{full adders}, test it on all 256 inputs, and give its size and delay.
\end{exercise}

\section{Problem: An arithmetic unit from one kind of gate}\label{sec:b1:combinational-circuits:problem}

\begin{problem}[{Weekend problem --- An arithmetic unit from one kind of gate}]\label{pb:b1:combinational-circuits:1}
A factory makes only NAND gates. How many does a half adder need? The program tries every circuit of $k$ two-input NAND
gates, for $k = 1, 2, \dots$; each signal is kept as its truth table on the four rows, a number of four bits.
\omcode{code/bachelor-1/27-combinational-circuits/listings/min_nand.py}{10}{27}{Every NAND circuit of $k$
gates}\label{lst:b1:combinational-circuits:min-nand}
\omoutput{code/bachelor-1/27-combinational-circuits/outputs/min_nand.txt}{The number of circuits tried, and the first
half adder found; signals 0 and 1 are $a$ and $b$, signal $i + 2$ the output of gate $i$.}\label{out:b1:combinational-circuits:min-nand}

\medskip\noindent\textbf{Part I --- Functions and forms.}
\begin{enumerate}
\item Give the truth tables of the sum and the carry of a half adder.
\item Write their canonical sums of products.
\item Why are \code{0b0011} and \code{0b0101} the truth tables of $a$ and $b$?
\item How does \code{\textasciitilde(x & y) & 0b1111} compute a NAND on the four rows at once?
\item How many \omterm{def:b1:combinational-circuits:boolean-function}{Boolean functions} of two variables are there? How many can a single NAND gate of $a$ and $b$ compute?
\item Why is a circuit of gates with no loop a \omterm{def:b1:combinational-circuits:combinational}{combinational circuit}?
\end{enumerate}

\medskip\noindent\textbf{Part II --- Translations into NAND.}
\begin{enumerate}[resume]
\item Translate the carry $ab$ into NAND gates, and count them.
\item Translate the sum $a \oplus b$ with four NAND gates, and check it on the four rows.
\item Read the circuit found for $k = 5$ and explain each gate.
\item Why may the two outputs share a gate?
\item Translate a \omterm{def:b1:combinational-circuits:adder}{full adder} into NAND gates by the same method, and count them.
\item What happens to the delay when each gate is replaced by its NAND translation?
\end{enumerate}

\medskip\noindent\textbf{Part III --- Adders and delays.}
\begin{enumerate}[resume]
\item Read the sizes and delays of \cref{out:b1:combinational-circuits:adders}.
\item Why is the delay of the \omterm{def:b1:combinational-circuits:adder}{ripple-carry adder} linear and that of a lookahead block constant?
\item Why do 8-bit lookahead adders have delay 6, and 64-bit ones 34?
\item How many inputs does the largest AND gate of a lookahead block have? What does that cost in hardware?
\item How many cases does the exhaustive check of the 8-bit adders try, and why is that a proof for 8 bits?
\item Why can the 64-bit adders not be checked the same way?
\end{enumerate}

\medskip\noindent\textbf{Part IV --- The smallest half adder.}
\begin{enumerate}[resume]
\item Show that the number of circuits of $k$ gates is $\prod_{i=0}^{k-1} \binom{i + 3}{2}$, and check it against the
output.
\item Why may a gate take the same signal twice?
\item Why do circuits with unused gates not spoil the search?
\item Why does the search for $k = 4$ prove that four gates are not enough?
\item Could a gate with three inputs, or a circuit with loops, do better? Does the search say?
\item State the named result, \emph{the half adder needs exactly five NAND gates}.
\end{enumerate}
\end{problem}
